\begin{table}[H]
\centering
\caption{Nama Daerah Perbandingan Cluster Biplot vs Kategori IPM (UNDP) (2022)} \label{tab:crossname2022}
\begin{adjustbox}{max width=\textwidth}
\begin{tabular}{|p{4cm}|p{10cm}|}
\hline
\textbf{Kategori} & \textbf{Daftar Wilayah} \\
\hline
Kl.1 \& Sangat Tinggi & - \\
\hline
Kl.1 \& Tinggi & Kep. Seribu, Bogor, Ciamis, Kuningan, Cirebon, Sumedang, Purwakarta, Karawang, Kota Tasikmalaya, Kota Banjar, Tangerang, Kota Cilegon, Kota Serang \\
\hline
Kl.1 \& Sedang & Sukabumi, Cianjur, Garut, Tasikmalaya, Majalengka, Indramayu, Subang, Bandung Barat, Pandeglang, Lebak, Serang \\
\hline
Kl.1 \& Rendah & - \\
\hline
Kl.2 \& Sangat Tinggi & - \\
\hline
Kl.2 \& Tinggi & - \\
\hline
Kl.2 \& Sedang & Pangandaran \\
\hline
Kl.2 \& Rendah & - \\
\hline
Kl.3 \& Sangat Tinggi & Kota Jakarta Selatan, Kota Jakarta Timur, Kota Jakarta Pusat, Kota Jakarta Barat, Kota Jakarta Utara, Kota Bandung, Kota Bekasi, Kota Depok, Kota Tangerang Selatan \\
\hline
Kl.3 \& Tinggi & Bandung, Bekasi, Kota Bogor, Kota Sukabumi, Kota Cirebon, Kota Cimahi, Kota Tangerang \\
\hline
Kl.3 \& Sedang & - \\
\hline
Kl.3 \& Rendah & - \\
\hline
\end{tabular}
\end{adjustbox}
\end{table}